\documentclass[12pt]{article}
\usepackage[utf8]{inputenc}
\usepackage[T1]{fontenc}
\usepackage[margin=1in]{geometry}
\usepackage{setspace}
\usepackage{xcolor}
\usepackage[colorlinks=true, urlcolor=blue!60!black, linkcolor=blue!60!black]{hyperref}
\setlength{\parindent}{2em}
\setlength{\parskip}{1.25em}
\renewcommand{\baselinestretch}{1.0}
\begin{document}
\begin{center}
{\LARGE \textit{BMW of North America, Inc. v. Gore}, 517 U.S. 559 (1996)\par}
\vspace{0.8em}
{\large Case Brief\par}
\end{center}
\vspace{1em}

\section{Memory Jogger}

In \href{https://www.courtlistener.com/opinion/118026/bmw-of-north-america-inc-v-gore/}{\textit{BMW of North America, Inc. v. Gore}}, the Supreme Court held for the first time in modern jurisprudence that a state punitive damages award---here, \$2 million on \$4,000 in actual damages for failing to disclose a \$601.37 presale paint repair---was ``grossly excessive'' under the Fourteenth Amendment's Due Process Clause, establishing a three-guidepost framework (reprehensibility, ratio, and comparable sanctions) for reviewing punitive awards.

\section{Detailed Case Facts}

In January 1990, Dr. Ira Gore, Jr., purchased a new black BMW sports sedan for \$40,750.88 from an authorized BMW dealership in Birmingham, Alabama (\href{https://supreme.justia.com/cases/federal/us/517/559/}{517 U.S. at 563}). After driving the vehicle for approximately nine months without noticing any flaw in its appearance, Dr. Gore took the car to ``Slick Finish,'' an independent detailing shop, to make it look ``snazzier than it normally would appear.'' The proprietor, Mr. Slick, detected that the top, hood, trunk, and quarter panels had been repainted at BMW's vehicle preparation center in Brunswick, Georgia, presumably after acid-rain damage sustained during transit from Germany.

Dr. Gore sued BMW of North America, Inc. (BMW), the American distributor, along with the German manufacturer and the Birmingham dealership, in Alabama state court. Relying on Alabama's 1907 fraud statute, Ala. Code \S~6-5-102 (1993)---which defines fraud as ``[s]uppression of a material fact which the party is under an obligation to communicate''---Dr. Gore alleged that BMW's nondisclosure of the repainting constituted fraudulent suppression and sought \$500,000 in compensatory and punitive damages. Under Ala. Code \S\S~6-11-20 and 6-11-21 (1993), punitive damages were authorized for ``gross, oppressive or malicious'' fraud.

At trial, BMW admitted that since 1983 it had maintained a nationwide policy regarding vehicles damaged during manufacture or transit: if repair costs exceeded 3\% of the manufacturer's suggested retail price (MSRP), the car was placed in company service and later sold as used; if repair costs did not exceed 3\% of MSRP, the car was repaired and sold as new without disclosing the repair to the dealer or customer (\href{https://supreme.justia.com/cases/federal/us/517/559/}{517 U.S. at 563--564}). Because repainting Dr. Gore's sedan cost \$601.37---roughly 1.5\% of its MSRP---BMW made no disclosure.

To establish compensatory damages, Dr. Gore presented testimony from a former Birmingham BMW dealer estimating that a repainted BMW lost roughly 10\% of its retail value, or \$4,000. To support punitive damages, Dr. Gore introduced evidence that since 1983 BMW had sold 983 refinished cars nationwide as new without disclosing paint repairs exceeding \$300, including 14 in Alabama. Multiplying \$4,000 by approximately 1,000 cars, Dr. Gore's counsel asked the jury for \$4 million in punitive damages. BMW countered that its 3\% threshold complied with the disclosure statutes of roughly 25 states (none of which required disclosure below 3\% of MSRP at the time of sale) and that it had never previously been adjudged liable for punitive damages for the policy; indeed, in an earlier Alabama suit, \textit{Yates v. BMW of North America, Inc.}, 642 So. 2d 937 (Ala. 1993), a jury had awarded compensatory damages but zero punitive damages. Prior to the verdict in Dr. Gore's case, BMW stopped shipping refinished cars to Alabama and two other states, and five days after the verdict, BMW adopted a nationwide policy of full disclosure (\href{https://supreme.justia.com/cases/federal/us/517/559/}{517 U.S. at 565--566, 579 n. 31}).

Historically and politically, the case arrived at the height of the 1990s national tort-reform debate. Business groups, insurers, and manufacturers had pressed the Supreme Court to rein in large state-court punitive verdicts, especially in jurisdictions such as Alabama. In \textit{Pacific Mutual Life Insurance Co. v. Haslip}, 499 U.S. 1 (1991), and \textit{TXO Production Corp. v. Alliance Resources Corp.}, 509 U.S. 443 (1993), the Court had acknowledged in theory that the Fourteenth Amendment Due Process Clause limits ``grossly excessive'' punitive awards, yet had upheld both awards. \textit{Gore} presented the vehicle through which a five-Justice majority finally gave substantive teeth to that constitutional boundary.

\section{Procedural History}

\begin{itemize}
\item \textbf{Trial Court (Circuit Court of Jefferson County, Alabama):} On June 12, 1992, the jury returned a verdict awarding Dr. Gore \$4,000 in compensatory damages and \$4 million in punitive damages against BMW (and also finding the dealership liable for compensatory damages and the German manufacturer liable for both) (\href{https://www.oyez.org/cases/1995/94-896}{Oyez}; \href{https://supreme.justia.com/cases/federal/us/517/559/}{517 U.S. at 565, 611 n. 1}). BMW moved post-trial to set aside the punitive award as grossly excessive under the Fourteenth Amendment Due Process Clause; the trial judge denied the motion.
\item \textbf{Alabama Supreme Court:} In \textit{BMW of North America, Inc. v. Gore}, 646 So. 2d 619 (Ala. 1994), the Alabama Supreme Court reversed the judgment against the German manufacturer for lack of personal jurisdiction, upheld the jury's finding of fraud, and applied the seven-factor excessiveness framework from \textit{Green Oil Co. v. Hornsby}, 539 So. 2d 218 (Ala. 1989). Crucially, the court held that the jury had improperly computed the \$4 million punitive award by multiplying \$4,000 by BMW's out-of-state sales without proof that those sales were unlawful where they occurred. Disclaiming any reliance on out-of-state acts, the court ordered a remittitur reducing the punitive damages award from \$4 million to \$2 million, deeming \$2 million ``constitutionally reasonable'' (\href{https://supreme.justia.com/cases/federal/us/517/559/}{517 U.S. at 566--567}).
\item \textbf{U.S. Supreme Court:} The Court granted certiorari, 513 U.S. 1125 (1995), to illuminate the standard identifying unconstitutionally excessive punitive damages awards, heard oral argument on October 11, 1995, and decided the case on May 20, 1996. By a 5--4 vote, the Court \textbf{reversed and remanded} to the Alabama Supreme Court to determine whether a new trial or a further independent remittitur was appropriate (\href{https://www.law.cornell.edu/supct/html/94-896.ZS.html}{Cornell LII Syllabus}).
\end{itemize}

\section{Judicial Votes}

The Supreme Court decided the case by a 5--4 margin, crossing conventional ideological lines:

\begin{itemize}
\item \textbf{Majority (5):} Justice Stevens (author), joined by Justices O'Connor, Kennedy, Souter, and Breyer.
\item \textbf{Concurrence (3):} Justice Breyer (author), joined by Justices O'Connor and Souter.
\item \textbf{Dissent (2):} Justice Scalia (author), joined by Justice Thomas.
\item \textbf{Dissent (2):} Justice Ginsburg (author), joined by Chief Justice Rehnquist.
\end{itemize}

\section{Holding}

\begin{itemize}
\item \textbf{Judgment:} Reversed and remanded. The \$2 million punitive damages award against BMW was ``grossly excessive'' and therefore violated the Due Process Clause of the Fourteenth Amendment.
\item \textbf{Rule Announced:} First, principles of state sovereignty and comity prohibit a State from imposing punitive damages to punish or deter a defendant's lawful conduct in other jurisdictions; an award must be justified solely by the State's interest in protecting its own consumers and economy. Second, elementary notions of constitutional fairness require that a defendant receive fair notice of both the conduct subject to punishment and the severity of the potential penalty. Courts evaluate whether a punitive award is unconstitutionally excessive using three guideposts: (1) the degree of reprehensibility of the defendant's conduct; (2) the disparity (ratio) between the actual or potential harm suffered by the plaintiff and the punitive damages award; and (3) the difference between the punitive remedy and the civil or criminal penalties authorized or imposed in comparable cases (\href{https://supreme.justia.com/cases/federal/us/517/559/}{517 U.S. at 572--575}).
\item \textbf{Holding vs. Dicta:} The narrow holding is that Alabama's \$2 million punitive award (a 500:1 ratio to \$4,000 in purely economic harm, absent bad faith or comparable statutory sanctions) violated due process. Several influential passages are technically dicta: (a) the majority's extended federalism discussion forbidding States from punishing lawful out-of-state conduct (which Justice Scalia and Justice Ginsburg pointed out was unnecessary because the Alabama Supreme Court had already disclaimed reliance on out-of-state sales in remitting the award to \$2 million); (b) the Court's suggestion in footnote 21 that out-of-state transactions remain admissible as evidence of reprehensibility; and (c) the majority's commentary that low compensatory awards or hard-to-detect injuries might justify higher ratios, while refusing to establish a rigid mathematical bright line.
\end{itemize}

\section{Analysis of Opinions}

\subsection{Majority Opinion (Justice Stevens)}

Justice Stevens begins by framing the inquiry around Alabama's legitimate state interests in punishment and deterrence (\href{https://supreme.justia.com/cases/federal/us/517/559/}{517 U.S. at 568}). Surveying the patchwork of automobile-repair disclosure laws across the 50 states---ranging from 3\% to 6\% of MSRP thresholds, dollar cutoffs, or common-law fraud rules---Stevens emphasizes that reasonable lawmakers disagree on optimal disclosure rules. Invoking principles of state sovereignty, comity, and the Dormant Commerce Clause (\textit{Healy v. Beer Institute}, 491 U.S. 324 (1989); \textit{Bonaparte v. Tax Court}, 104 U.S. 592 (1881)), Stevens concludes that Alabama may protect its own consumers, but may not use punitive damages to punish BMW for conduct that was lawful in other states or to force a nationwide policy change. Thus, the \$2 million award had to be judged exclusively against BMW's 14 Alabama sales and the interests of Alabama consumers.

Turning to the excessiveness analysis, Stevens grounds constitutional review in ``fair notice'' and applies the three guideposts:
\begin{enumerate}
\item \textbf{Degree of Reprehensibility:} Calling reprehensibility ``[p]erhaps the most important indicium'' (\href{https://supreme.justia.com/cases/federal/us/517/559/}{517 U.S. at 575}), Stevens finds none of the aggravating factors ordinarily warranting a severe sanction. The harm was purely economic; the repainting did not affect performance, safety, or even appearance for nine months; and BMW showed no reckless disregard for health or safety. Nor could BMW be labeled a knowing recidivist, because a corporate executive could reasonably read the 25 state disclosure statutes (including thresholds of 3\% or higher) as safe harbors, and BMW promptly changed its policy once adjudged unlawful. There were no deliberate false statements, affirmative misconduct, or concealment of improper motive.
\item \textbf{Ratio:} Comparing the \$2 million punitive award to Dr. Gore's \$4,000 actual loss yields a ``breathtaking 500 to 1'' ratio---and even aggregating all 14 Alabama buyers yields \$56,000, making the award 35 times total statewide harm (\href{https://supreme.justia.com/cases/federal/us/517/559/}{517 U.S. at 582 \& n. 35}). Unlike \textit{Haslip} (a ratio slightly over 4:1) or \textit{TXO} (at most 10:1 when counting potential harm), BMW's policy threatened no additional potential harm. While rejecting a fixed mathematical formula, Stevens concludes that a 500:1 ratio in a purely economic case surely ``raise[s] a suspicious judicial eyebrow.''
\item \textbf{Sanctions for Comparable Misconduct:} Giving substantial deference to legislative judgments, Stevens compares the \$2 million award to Alabama's maximum civil penalty of \$2,000 under its Deceptive Trade Practices Act (Ala. Code \S~8-19-11(b)) and comparable state fines ranging from \$50 to \$10,000 (\href{https://supreme.justia.com/cases/federal/us/517/559/}{517 U.S. at 583--584}). Neither state statutes nor judicial precedents gave BMW fair notice that a first violation (or 14 sales) could trigger a multimillion-dollar sanction, nor was there any basis to assume a far more modest penalty would not have deterred future nondisclosure in Alabama.
\end{enumerate}

\subsection{Concurring Opinion (Justice Breyer)}

Justice Breyer, joined by Justices O'Connor and Souter, writes separately to explain why the ``strong presumption of validity'' ordinarily accorded to state judgments produced by fair procedures (\textit{TXO}, 509 U.S. at 457) is overcome here (\href{https://supreme.justia.com/cases/federal/us/517/559/}{517 U.S. at 586--597}). Breyer grounds his analysis in rule-of-law principles tracing back to Magna Carta: due process protects against arbitrary coercion by requiring legal standards that meaningfully constrain decisionmakers' discretion.

Examining Alabama's punitive damages regime in operation, Breyer identifies five structural deficiencies: (1) Alabama's punitive damages statute defines ``fraud,'' ``malice,'' and ``oppression'' so broadly that it fails to distinguish minor omissions from egregious misconduct; (2) the Alabama Supreme Court's application of the seven \textit{Green Oil} factors stripped them of constraining force (for instance, finding a ``reasonable relationship'' between \$56,000 in statewide harm and a \$2 million award, and relying on defendant's wealth to inflate rather than cabin the verdict); (3) Alabama courts articulated no constraining economic deterrence theory (such as multiplying total harm by the inverse probability of detection) to explain the \$2 million figure; (4) historical English and American exemplary awards, properly adjusted for inflation in an appendix, offered no precedent for such a massive award for minor economic fraud; and (5) Alabama had no legislative caps cabining jury discretion. Combining this absence of constraining standards with the severe disproportionality of the award overcame the presumption of validity.

\subsection{Dissenting Opinion (Justice Scalia)}

Justice Scalia, joined by Justice Thomas, delivers a sharp originalist and textualist dissent (\href{https://supreme.justia.com/cases/federal/us/517/559/}{517 U.S. at 598--607}). He rejects substantive due process review of punitive damages altogether, arguing that the Fourteenth Amendment's Due Process Clause guarantees only \textit{procedural} fairness---namely, an opportunity to contest the reasonableness of a jury's damages award in state court---and is not ``a secret repository of substantive guarantees against `unfairness.'\,'' At the time the Fourteenth Amendment was ratified, punitive damages were understood as the community's expression of indignation assessed by a jury; the Lochner-era penalty cases cited by the majority (such as \textit{Seaboard Air Line Railway Co. v. Seegers}, 207 U.S. 73 (1907)) were shallowly rooted and never invalidated a civil damages verdict as merely ``too big.''

Scalia further argues that the majority's extraterritoriality discussion is ``the purest dicta'' because the Alabama Supreme Court already disclaimed reliance on out-of-state sales, and he contends that the three guideposts ``mark a road to nowhere.'' Calling them ``crisscrossing platitudes'' that leave open unnamed additional factors and a deterrence loophole, Scalia warns that the Court has federalized state tort law without giving lower courts workable rules. Logically, he adds, elevating ``fairness'' in punitive awards to substantive due process implies that any allegedly unreasonable compensatory award or evidentiary-sufficiency dispute in state civil litigation could become a federal constitutional issue.

\subsection{Dissenting Opinion (Justice Ginsburg)}

Justice Ginsburg, joined by Chief Justice Rehnquist, dissents on institutional federalism and judicial-restraint grounds (\href{https://supreme.justia.com/cases/federal/us/517/559/}{517 U.S. at 607--619}). She emphasizes that the Alabama Supreme Court had already cured the jury's improper nationwide multiplier (to which BMW's trial counsel had failed to object) and had independently reviewed the award under the \textit{Green Oil} factors previously approved in \textit{Haslip}.

Ginsburg warns that the Court is unwisely entering a domain traditionally reserved to the States without any bright-line test---leaving only a vague ``raised eyebrow'' standard under which ``[t]oo big is, in the end, the amount at which five Members of the Court bridle.'' Because state-court punitive verdicts come directly to the Supreme Court on certiorari rather than through federal habeas review, the Supreme Court will stand alone as the only federal court policing state punitive awards. Appending a detailed survey of state legislative reforms---including statutory caps, split-recovery allocations to state funds, and mandatory bifurcation---Ginsburg argues that active state courts and legislatures should be left to address punitive-damages reform themselves.

\section{Examples: Future Applications}

\begin{enumerate}
\item \textbf{Same Side (Unconstitutionally Excessive):} A national insurer maintains an internal claims-handling practice that systematically underpays policyholders' minor property-damage claims by 5\% under ambiguous policy language that several states' insurance commissioners have not prohibited. A policyholder suffers \$10,000 in actual economic loss and wins \$15 million in punitive damages (a 1,500:1 ratio) based largely on evidence of the insurer's nationwide claims volume, while state law sets a \$5,000 administrative fine per violation. Under \textit{Gore}, this award violates due process: the harm is purely economic, the ratio is extreme, out-of-state lawful practices cannot be punished, and comparable statutory fines are tiny.
\item \textbf{Same Side (Unconstitutionally Excessive):} A software vendor sells a consumer productivity suite without disclosing that a minor background telemetry feature reduces battery life by 3\%. A state jury awards a single buyer \$200 in compensatory damages and \$500,000 in punitive damages (a 2,500:1 ratio), where no physical safety risk exists and state consumer-protection statutes cap civil penalties at \$2,500. Applying \textit{Gore}'s three guideposts---low reprehensibility, staggering ratio, and modest statutory penalties---the award is grossly excessive.
\item \textbf{Opposite Side (Constitutionally Permissible):} A pharmaceutical manufacturer deliberately conceals internal clinical-trial data showing that its over-the-counter painkiller causes fatal liver failure, affirmatively falsifies FDA submissions, and continues selling the drug for years to maximize profits. A surviving spouse recovers \$1 million in compensatory damages and \$4 million in punitive damages (a 4:1 ratio), and criminal fraud statutes authorize multimillion-dollar corporate fines and imprisonment. Under \textit{Gore}, the award easily survives due process: reprehensibility is extreme (physical harm, reckless disregard for health and safety, affirmative deceit), the 4:1 ratio falls well within \textit{Haslip} bounds, and comparable criminal sanctions are severe.
\item \textbf{Opposite Side (Constitutionally Permissible):} A landlord intentionally disconnects heat and water in mid-winter to force an elderly, low-income tenant out of a rent-controlled apartment, causing \$5,000 in documented economic and relocation losses. A jury awards \$35,000 in punitive damages (a 7:1 ratio), where municipal housing codes authorize up to \$25,000 in civil fines and misdemeanor jail time for intentional utility shutoffs. Because the conduct targets a financially vulnerable victim, threatens physical health, reflects intentional malice, bears a modest single-digit ratio, and aligns with statutory penalties, it passes \textit{Gore}.
\item \textbf{Fence-Sitter (Borderline Case):} A corporate competitor intentionally trespasses onto a small business's property to dump non-toxic construction debris and sabotage a bid, causing only \$500 in actual property-cleanup costs because the owner discovers and removes the debris immediately, though the sabotage could have cost the business a \$200,000 contract. A jury awards \$150,000 in punitive damages (a 300:1 ratio to actual damages, but less than 1:1 to potential harm), while state criminal mischief statutes carry a \$5,000 maximum fine. This case sits squarely on \textit{Gore}'s fence: the raw actual-damages ratio and modest statutory fine point toward excessiveness, yet \textit{Gore} explicitly notes that intentional misconduct causing small actual damages---or significant \textit{potential} harm, as in \textit{TXO}---can justify a much higher ratio.
\end{enumerate}

\section{Critique}

\textit{BMW of North America, Inc. v. Gore} has drawn sustained scholarly scrutiny from both progressive and conservative/originalist camps (\href{https://www.everycrsreport.com/reports/RL33773.html}{CRS Report RL33773}; \href{https://scholarship.law.vanderbilt.edu/cgi/viewcontent.cgi?article=1305&context=vlr}{Vanderbilt Law Review}).

\textbf{Originalist and Textualist Critique:} Echoing Justice Scalia's dissent, originalists and textualists argue that \textit{Gore} resurrects economic substantive due process reminiscent of \textit{Lochner v. New York}, 198 U.S. 45 (1905), without textual or historical mooring in the Fourteenth Amendment. At ratification in 1868, ``due process of law'' guaranteed established judicial procedures, not federal appellate review of jury-assessed tort remedies. Critics also highlight the irony of the Court's federalism rhetoric: to prevent one state from regulating nationwide conduct, the Court federalized traditional state tort law and displaced state supreme courts and legislatures that were actively enacting their own tort-reform caps and bifurcation rules.

\textbf{Progressive and Institutional Critique:} Progressive scholars and trial-law advocates criticize \textit{Gore} from the opposite vantage point. First, they contend that anchoring punitive awards to single-plaintiff compensatory ratios undermines optimal deterrence when a wealthy corporation commits widespread, low-dollar consumer fraud: if only a small fraction of cheated consumers ever discover the fraud and sue, limiting punitive damages to a small multiple of a single buyer's \$4,000 loss makes fraud systematically profitable. Second, critics argue that the Court's ``fair notice'' rationale rings hollow when applied to sophisticated multinational corporations advised by counsel, while shielding corporate wealth from community accountability. Third, scholars note the doctrinal instability of the three guideposts: because \textit{Gore} refused to set a bright-line ratio, lower courts struggled with unpredictable outcomes until \textit{State Farm Mutual Automobile Insurance Co. v. Campbell}, 538 U.S. 408 (2003), tightened \textit{Gore} into a presumptive single-digit ratio ceiling---a move that critics argue further transformed a flexible due-process floor into a rigid, judge-made tort-reform cap.

\textbf{Logical Weaknesses and Competing Values:} The opinion sits atop an unresolved tension between two competing constitutional values: protecting defendants from arbitrary, disproportionate state coercion versus respecting democratic federalism and the common-law jury's role in expressing moral condemnation. Moreover, as Justice Scalia noted, the majority's insistence that out-of-state lawful conduct cannot be punished sits uneasily beside footnote 21's permission to admit out-of-state conduct as evidence of ``reprehensibility,'' inviting juries to do indirectly what the Court barred them from doing directly.

\section{Key Quotations}

\begin{quote}
``Punitive damages may properly be imposed to further a State's legitimate interests in punishing unlawful conduct and deterring its repetition. ... Only when an award can fairly be categorized as `grossly excessive' in relation to these interests does it enter the zone of arbitrariness that violates the Due Process Clause of the Fourteenth Amendment.'' \\
--- Justice Stevens, Opinion of the Court (\href{https://supreme.justia.com/cases/federal/us/517/559/}{517 U.S. at 568}).
\end{quote}

\begin{quote}
``Elementary notions of fairness enshrined in our constitutional jurisprudence dictate that a person receive fair notice not only of the conduct that will subject him to punishment, but also of the severity of the penalty that a State may impose. Three guideposts, each of which indicates that BMW did not receive adequate notice of the magnitude of the sanction that Alabama might impose for adhering to the nondisclosure policy adopted in 1983, lead us to the conclusion that the \$2 million award against BMW is grossly excessive: the degree of reprehensibility of the nondisclosure; the disparity between the harm or potential harm suffered by Dr. Gore and his punitive damages award; and the difference between this remedy and the civil penalties authorized or imposed in comparable cases.'' \\
--- Justice Stevens, Opinion of the Court (\href{https://supreme.justia.com/cases/federal/us/517/559/}{517 U.S. at 574--575}).
\end{quote}

\begin{quote}
``In most cases, the ratio will be within a constitutionally acceptable range, and remittitur will not be justified on this basis. When the ratio is a breathtaking 500 to 1, however, the award must surely `raise a suspicious judicial eyebrow.'\,'' \\
--- Justice Stevens, Opinion of the Court (\href{https://supreme.justia.com/cases/federal/us/517/559/}{517 U.S. at 583}).
\end{quote}

\begin{quote}
``I do not regard the Fourteenth Amendment's Due Process Clause as a secret repository of substantive guarantees against `unfairness'---neither the unfairness of an excessive civil compensatory award, nor the unfairness of an `unreasonable' punitive award. What the Fourteenth Amendment's procedural guarantee assures is an opportunity to contest the reasonableness of a damages judgment in state court; but there is no federal guarantee a damages award actually be reasonable.'' \\
--- Justice Scalia, Dissenting (\href{https://supreme.justia.com/cases/federal/us/517/559/}{517 U.S. at 598--599}).
\end{quote}

\begin{quote}
``What is the Court's measure of too big? Not a cap of the kind a legislature could order, or a mathematical test this Court can divine and impose. Too big is, in the end, the amount at which five Members of the Court bridle.'' \\
--- Justice Ginsburg, Dissenting (\href{https://supreme.justia.com/cases/federal/us/517/559/}{517 U.S. at 613 n. 5}).
\end{quote}

\end{document}
